% ===== 原文 pp.9–13（PDF p.19–23） =====
% [原文 p.9]
\chapter{四维规范理论的对偶}

我们先回顾四维（4d）中的 $\cN=4$ 超杨-米尔斯（super Yang-Mills，SYM）理论，以及它的对偶性质与把该理论扭转（twisting）后所能得到的各种拓扑场论（TFT）。

\section{电磁对偶的非阿贝尔推广}

在 Yang-Mills 理论出现之后，人们自然会问：电磁对偶的某种版本能否推广到规范群 $G$ 非阿贝尔的理论。的确，既然对偶会反转耦合参数 $e^2$，对偶就有可能为非阿贝尔规范理论的强耦合行为提供一个强有力的窗口：只需在对偶的、弱耦合的那一侧研究强耦合物理的某个问题即可。

% [Tutor Note: p.9 术语首现。Wilson line operator（Wilson 线算子）与 t'Hooft disorder line operator（t'Hooft 无序线算子）是电磁对偶两端的“电性”与“磁性”对象。原文通篇把这名字印作 “t'Hooft”（撇号在 t 之后），国际通行写法是 “'t Hooft”（撇号在 t 之前，指 G. 't Hooft 姓氏的冠词）；按“照原文译、不改写原文”的原则，本译本正文一律保留原文的 “t'Hooft”，仅在此说明通行写法。]

这样的非阿贝尔对偶如果存在，就必须把一个 Wilson 线算子（“电性源”的非阿贝尔推广）与一个 t'Hooft 无序线算子（“磁性源”的非阿贝尔推广）互相交换。但在考察非阿贝尔规范理论中 t'Hooft 算子的分类时，Goddard、Nuyts 与 Olive [25] 发现了故事中的一个关键转折：Wilson 算子由群 $G$ 的不可约表示标记，而 t'Hooft 算子却被 $G$ 的余权格的某个基本域中的元素标记，这些元素反过来又标记另一个李群 $\LG$ 的不可约表示。

% [Tutor Note: p.9 可验证的易错点。“余权格的基本域”指 Weyl 群在余权格上的基本域，其元素与 $\LG$ 的支配权一一对应，因此“由 $\LG$ 的不可约表示标记”与“由余权格基本域的元素标记”是同一件事的两种说法。核对：$G=SU(N)$ 时 $\LG=SU(N)/\Z_N$，余权格与权格重合，基本域的元素就是支配权，对应 $SU(N)$ 不可约表示的最高权。]

正如我们在上一章所强调的，由交换余权格与权格相联系的李群对 $G$ 与 $\LG$ 在数论中的朗兰兹纲领里起核心作用；正因为如此，$\LG$ 被称为 $G$ 的朗兰兹对偶群。

这些考虑启发了 Montonen 与 Olive [23] 去猜想：在量子层面，规范群为 $G$ 的 Yang-Mills 理论与规范群为 $\LG$、且耦合反转的 Yang-Mills 理论对偶。

% [原文 p.10]
然而，人们很快就意识到这个朴素的对应不可能完全正确；首先，耦合之间的这种简单对应在重整化下不会保持。Osborn [24] 在 Witten 与 Olive [26] 的工作基础上把猜想改进为：Yang-Mills 理论的超对称版本——即规范群为 $G$ 的 $\cN=4$ 超 YM 理论——与规范群为 $\LG$、耦合反转的 $\cN=4$ 超 YM 理论对偶。由于 $\cN=4$ 是共形场论，耦合之间的反转对应不会被重整化破坏。

尽管严格说来，$\cN=4$ 的对偶在数学严格性的层面仍未获证明，物理学家在过去三十年中已经积累了支持这一猜想的丰富证据；时至今日它已被视为既定事实，本文也将如此对待。

下面我们更详细地回顾 $\cN=4$ 理论及其对偶性质。

\section{\texorpdfstring{$\cN=4$、$d=4$、规范群 $G$}{N=4、d=4、规范群 G}}

$\cN=4$ 理论是四维 Yang-Mills 理论的最大超对称延拓。它由作用在时空——一个四维流形 $M$——上的作用量来定义。为了写下作用量，让我们先固定一些规范理论的约定。本章余下部分将紧密遵循 [1] 中的记号约定与论述。

\subsection{规范理论约定}

设 $G$ 是连通约化李群，李代数为 $\fg$。我们一次性固定 $G$ 的一个忠实的 $N$ 维“定义”表示，从而把 $\fg$ 等同于一族 $N\times N$ 反厄米矩阵。设 $T_a$（$a=1,\dots,\dim G$）为 $\fg$ 的一组基，并用它定义（实）结构常数 $f^{c}{}_{ab}$：
\[ [T_a,T_b] = f^{c}{}_{ab}T_c \]

设 $P$ 是 $M$ 上的主 $G$-丛，平凡化邻域为 $U^{(\alpha)}\subset M$，局部截面为 $s^{(\alpha)}:U^{(\alpha)}\to P$。规范场将是 $P$ 上的一个联络 $A$，我们在局部把它表示为一个 $\fg$-值 1-形式
\[ A^{(\alpha)} \equiv s^{(\alpha)*}A = A^{(\alpha)}_\mu \dd x^\mu = A^{(\alpha)a}_\mu T_a \dd x^\mu \in \Omega^1(U^{(\alpha)},\fg) \]
其曲率由 $F^{(\alpha)} = dA^{(\alpha)} + \frac{1}{2}[A^{(\alpha)},A^{(\alpha)}]$ 给出。今后我们往往省略局部定义量上的标号 $(\alpha)$。

我们采用记号
\[ \eta \in \Omega^k(M, ad\,P) \]
来表示 $\eta$ 是 $M$ 上一个取值于 $ad\,P$ 的 $k$-形式，而 $ad\,P$ 是由 $P$ 经 $G$ 的
% [原文 p.11]
伴随表示得到的相伴向量丛。$\eta$ 关于 $A$ 的协变导数由下式给出：
\[ d_A\eta = d\eta + [A,\eta] \]

参数为 $\epsilon\in\Omega^0(M,ad\,P)$ 的 $P$ 的无穷小自同构把 $A$ 与 $\eta$ 变换为
\[ \delta_\epsilon A = -d_A\epsilon \]
\[ \delta_\epsilon\eta = [\epsilon,\eta] \]

% [Tutor Note: p.11 补出省略论证。$\delta_\epsilon A=-d_A\epsilon$ 来自规范变换 $A\mapsto g^{-1}Ag+g^{-1}dg$ 在 $g=\ee^{\epsilon}\simeq1+\epsilon$ 时的展开：$\delta A=d\epsilon+[\epsilon,A]=-d_A\epsilon$（注意原文的 $A$ 反厄米，故 $g^{-1}dg\simeq d\epsilon$ 而 $[\epsilon,A]$ 的符号与厄米约定相反）；$\delta_\epsilon\eta=[\epsilon,\eta]$ 只是 $ad\,P$ 的截面按伴随表示变换。自洽核对：由这两条可推出 $\delta_\epsilon F=[\epsilon,F]$，即曲率按伴随表示变换。]

设“$\Tr$”为 $\fg$ 上一个 Ad-不变的负定度规，其归一化使得
\[ \frac{1}{8\pi^2}\int_M \Tr F\wedge F \]
当 $M=S^4$ 时取整数值。当需要我们直接提到这个度规时，记它为 $\kappa_{ab}$，其逆记为 $\kappa^{ab}$。

% [Tutor Note: p.11 归一化核对。$\frac{1}{8\pi^2}\int_M\Tr F\wedge F$ 取整数值与 $\kappa_{ab}=\Tr(T_aT_b)$ 负定是同一件事的两面：对反厄米生成元 $\Tr(T_aT_b)<0$。以 $SU(2)$ 为例，$-\frac{1}{8\pi^2}\Tr F\wedge F=c_2$，其积分给出瞬子数 $k\in\Z$，故 (2.2) 中系数 $\theta/8\pi^2$ 正是让 $\theta$ 成为以 $2\pi$ 为周期的角度参数的标准归一化（归一化若差一个整数倍，$\theta\to\theta+2\pi$ 就不再是理论的对称性）。]

\subsection{作用量与超对称}

起初我们取 $M$ 为平坦的 $\R^{1,3}$，采用 Lorentz 号差。除规范场 $A$（主丛 $P$ 上的联络）之外，物质内容由六个 $ad\,P$-值标量场
\[ \phi_1,\dots,\phi_6 \in \Omega^0(M, ad\,P) \]
以及八个费米子场
\begin{gather*}
\lambda^1,\dots,\lambda^4 \in \Omega^0(M, ad\,P\otimes\cS_+) \\
\bar\lambda_1,\dots,\bar\lambda_4 \in \Omega^0(M, ad\,P\otimes\cS_-)
\end{gather*}
组成，其中 $\cS_\pm$ 是 $M$ 上的两个旋量丛。$\cN=4$ 超 YM 的作用量由下式给出
\begin{equation}
S = \frac{1}{e^2}\int_M \dd^4x\,\Tr\left\{\frac{1}{2}F_{\mu\nu}F^{\mu\nu} + D_\mu\phi_i D^\mu\phi^i + \frac{1}{2}[\phi_i,\phi_j][\phi^i,\phi^j]\right\} + \ldots \tag{2.1}
\end{equation}
其中 $\ldots$ 表示作用量中依赖于费米子的那些项（此处略去）。此外，我们还加入一个测量丛 $P$ 的第二陈类的拓扑项：
\begin{equation}
S_{top} = -\frac{\theta}{8\pi^2}\int_M \Tr F\wedge F \tag{2.2}
\end{equation}

由于该理论在 $\theta\to\theta+2\pi$ 下显然不变（路径积分权重不变），把 $\theta$ 视为一个角度参数是很自然的。参数 $e$ 与 $\theta$ 可以方便地
% [原文 p.12]
合并成一个复化耦合参数：
\[ \tau = \frac{\theta}{2\pi} + \frac{4\pi\ii}{e^2} \]

% [Tutor Note: p.11 文字层陷阱 + p.12 原文笔误。(2.1) 中第三项原文写作 $D_\mu\phi_iD^\mu\phi^i$（第二个 $\phi$ 带上指标），文字层抽取成 $D_\mu\phi_iD^\mu\phi_i$ 是错的，须以图像为准；用 $\Tr$ 的 Ad-不变性可验证该项与 $F^2$ 项一样规范不变。另：p.12 第 8 行 “with respect to which which the six $\phi$s are rotated” 重复了 which，正确形式为 “with respect to which the six $\phi$s...”；同页 §2.3 首段的 “at the the quantum level” 亦重复 the。]

除 $Spin(1,3)$ Lorentz 对称性外，该理论还在一个 $Spin(6)$ 对称性下不变：在这个对称性下六个 $\phi$ 作为一个 $\mathbf{6}$ 相互旋转，而 $\lambda$ 作为 $\mathbf{4}$、$\bar\lambda$ 作为 $\bar{\mathbf{4}}$ 变换。我们称这个群为该理论的 R-对称群，记作 $Spin(6)_R\simeq SU(4)_R$。

此外，该理论在费米型超对称 $Q^1_{\dot A},\dots,Q^4_{\dot A}$ 下不变，它们在玻色对称群
\[ Spin(1,3)\times SU(4)_R \simeq SL(2,\C)\times SU(4)_R \]
下作为 $(\mathbf{2},\mathbf{4})$ 变换；而 $\bar Q_{A1},\dots,\bar Q_{A4}$ 则作为 $(\bar{\mathbf{2}},\bar{\mathbf{4}})$ 变换。这里 $A=1,2$ 是 $SL(2,\C)$ 的 $\mathbf{2}$ 指标，$\dot A=1,2$ 是 $SL(2,\C)$ 的 $\bar{\mathbf{2}}$ 指标。

\section{S-对偶}

让我们考察 $\cN=4$ 的对偶性质。首先，把复化耦合 $\tau$ 平移 $\tau\to\tau+1$ 对应于角度 $\theta$ 的 $2\pi$ 平移，这显然是理论的一个对称性。

不那么平凡的是，把 Montonen-Olive 猜想推广到 $\cN=4$ SYM 后断言：在量子层面，规范群为单链群 $G$、耦合为 $\tau$ 的理论，与把 $G$ 换成其朗兰兹对偶群 $\LG$、把电性源与磁性源互换、并把 $\tau$ 换成
\[ {}^L\tau = -\frac{1}{\tau} \]
后得到的理论同构。这个 S-对偶与 $\theta$ 角的平移结合起来，生成一个 $SL(2,\Z)$ 对偶群。

当 $G$ 非单链时，耦合在 S-对偶下的变换为
\[ {}^L\tau = -\frac{1}{n_{\fg}\tau} \]
其中 $n_{\fg}=1,2,3$ 是该群的链数，即长根与短根长度平方之比（根位于 $\mathfrak{t}^*$，即 $\fg$ 的 Cartan 子代数的对偶，度规取 Ad-不变的 $\Tr$）。当 $n_{\fg}=2$ 或 $3$ 时，$\tau$ 上的 S-对偶映射与 $\theta$ 角的平移结合起来，生成 $SL(2,\R)$ 的一个离散子群，称为 Hecke 群。关于朗兰兹对偶李群对的更多细节与实例，见 [2] 的附录 A。

正如 [1] 中所强调的，联系对偶理论中耦合强度的公式必须谨慎解释。
% [原文 p.13]
大多数单李群 $G$ 与其对偶有相同的李代数，即对大多数单群有 $\fg=\Lg$。这时在写下两个理论的作用量时，把两个向量空间 $\fg$ 与 $\Lg$ 等同起来是自然的；的确，检查上面的作用量即可看出，其形式只依赖于李代数的数据。这样一来，人们就可以比较对偶理论中的耦合。另一方面，李群对
\[ G = Spin(2k+1) \quad\longleftrightarrow\quad {}^LG = Sp(k)/\Z_2 \]
的李代数并不同构：$\fg\neq\Lg$。因此，为了比较耦合必须做一个非标准的归一化选择；与 [1] 一样，我们采用如下约定：上面关于 ${}^L\tau$ 的公式继续成立，并取 $n_{\fg}=2$。

% [Tutor Note: p.13 可验证的易错点。$Sp(k)/\Z_2$ 与 $Spin(2k+1)$ 的维数同为 $k(2k+1)$，但李代数分别是 $C_k$ 与 $B_k$，并不同构，因此不能把两边的根长归一化等同起来。核对原文的取法：$B_k$ 与 $C_k$ 的链数都等于 $2$（长根与短根长度平方之比为 $2$），所以“取 $n_{\fg}=2$”并非任意约定，而是与链数定义一致的；反过来说，若这里误取 $n_{\fg}=1$，${}^L\tau=-1/\tau$ 与 $\theta$ 平移生成的群就不再同构于 $Hecke$ 群，与原文 p.12 的陈述矛盾。]

要知道 S-对偶如何变换 $\cN=4$ 的超对称，对我们来说将很重要。人们可以从超对称代数中心荷的变换推断出 [1]：左手超对称被乘上一个相位，右手超对称被乘上相反的相位，即
\[ \bar Q_{Ai} \longrightarrow \ee^{\ii\alpha/2}\bar Q_{Ai}, \qquad Q^i_{\dot A} \longrightarrow \ee^{-\ii\alpha/2}Q^i_{\dot A}, \qquad \text{其中}\quad \ee^{\ii\alpha}\equiv\frac{\abs{\tau}}{\tau} \]

% [Tutor Note: p.13 可验证的易错点。$\ee^{\ii\alpha}\equiv|\tau|/\tau$，在 $\theta=0$ 时 $\tau=4\pi\ii/e^2$ 纯虚，故 $\tau/|\tau|=\ii$ 而 $|\tau|/\tau=-\ii$（两者互为倒数，切勿混用）。自洽检验：S-对偶作用两次给出 $\tau\mapsto-1/\tau\mapsto\tau$，而由本节末的 $t\mapsto\frac{\tau}{|\tau|}t$（见 p.17）作用两次给出 $t\mapsto-t$；这与 §2.5.1 所说“理论在 $t$ 与 $-t$ 处等价”（由 $R$-对称群中心的作用）恰好吻合，说明 $Q$ 的相位约定与 $t$ 的变换规则是自洽的。]

\section{扭转}

我们的主要兴趣并不是 $\cN=4$ 理论本身，而是由 $\cN=4$ 经扭转（twisting）手续得到的一族拓扑场论。扭转（由 E. Witten 在 1980 年代关于四维流形 Donaldson 不变量的工作 [45] 中引入）是一种从给定的物理超对称场论出发构造拓扑场论的方法。

% [Tutor Note: p.13 术语首现。twisting 译为“扭转”，twisted theory 为“扭转后的理论”；文献 [45] 即 Witten 关于 Donaldson 不变量的原始工作，也是“拓扑扭转”这一构造的出处，与 §2.4.1 中的 Donaldson-Witten 扭转同名但这里是泛指扭转手续本身。]

我们扭转的动机是：利用扭转后理论关联函数的度规无关性，我们就能以最小的代价精确计算扭转理论中的某些量，而对未扭转的 $\cN=4$ 而言相应的分析会复杂得多。例如在第 6 章中，我们将利用度规无关性，通过紧化到三维并考察边界条件来直接计算面算子的谱。

一个相关的动机是在比平坦 $\R^4$ 更一般的四维流形上研究超对称场论的对偶。一般来说，超对称理论的费米型对称性无法在从 $\R^4$ 过渡到更一般的弯曲流形时存活下来（除非几何允许协变常数旋量这种受限情形）。但正如我们马上会看到的，扭转会产生至少一个费米型荷 $Q$，它在各坐标卡之间作为 0-形式变换，从而在把时空取为一般四维流形之后仍被保留。例如在后面的章节中，我们将研究形如 $\Sigma\times C$（$\Sigma$ 与 $C$ 为黎曼面）以及形如 $W\times S^1$（$W$ 为三维流形，$S^1$ 为圆周）的时空上的对偶。
